\ifdefined\gkver\else\def\gkver{student}\fi
\documentclass[\gkver]{gaokaozhenti}
\begin{document}

\chapter{2007年四川理综}
%% paperType: 理综物理部分

\begin{enumerate}
\item
%% number: 14
%% paperName: 2007年四川理综第 14 题 6 分
%% typeId: 1
%% type: 单选题
%% chapter: 气体性质
%% point: 热力学第一定律
%% method: 无
%% score: 6
%% degree: 650
%% duplicateId: 0
%% body:
如图所示，厚壁容器的一端通过胶塞插进一只灵敏温度计和一根气针，另一端有个用卡子卡住的可移动胶塞。用打气筒慢慢向筒内打气，使容器内的压强增加到一定程度，这时读出温度计示数。打开卡子，胶塞冲出容器后\xzanswer{C}

\onepicture[w=6.0cm]{figs/14.png}

\fourchoices[answer=C]
{温度计示数变大，实验表明气体对外界做功，内能减少}
{温度计示数变大，实验表明外界对气体做功，内能增加}
{温度计示数变小，实验表明气体对外界做功，内能减少}
{温度计示数变小，实验表明外界对气体做功，内能增加}

%% answer: C
%% memo:
\memoanswer{
本题原卷及材料中未提供详解，待补充。
}

\item
%% number: 15
%% paperName: 2007年四川理综第 15 题 6 分
%% typeId: 1
%% type: 单选题
%% chapter: 电磁感应
%% point: 切割磁感线电动势
%% method: 无
%% score: 6
%% degree: 450
%% duplicateId: 0
%% body:
如图所示，矩形线圈 $abcd$ 在匀强磁场中可以分别绕垂直于磁场方向的轴 $P_{1}$ 和 $P_{2}$ 以相同的角速度匀速转动，当线圈平面转到与磁场方向平行时\xzanswer{A}

\onepicture[w=4.0cm]{figs/15.png}

\fourchoices[answer=A]
{线圈绕 $P_{1}$ 转动时的电流等于绕 $P_{2}$ 转动时的电流}
{线圈绕 $P_{1}$ 转动时的电动势小于绕 $P_{2}$ 转动时的电动势}
{线圈绕 $P_{1}$ 和 $P_{2}$ 转动时电流的方向相同，都是 $a\to b\to c\to d$}
{线圈绕 $P_{1}$ 转动时 $dc$ 边受到的安培力大于绕 $P_{2}$ 转动时 $dc$ 边受到的安培力}

%% answer: A
%% memo:
\memoanswer{
本题原卷及材料中未提供详解，待补充。
}

\item
%% number: 16
%% paperName: 2007年四川理综第 16 题 6 分
%% typeId: 1
%% type: 单选题
%% chapter: 原子和原子核
%% point: 天然放射性
%% method: 无
%% score: 6
%% degree: 650
%% duplicateId: 0
%% body:
关于天然放射现象，下列说法正确的是\xzanswer{D}

\fourchoices[answer=D]
{放射性元素的原子核内的核子有半数发生变化所需的时间就是半衰期}
{放射性物质放出的射线中，$\alpha$ 粒子动能很大，因此贯穿物质的本领很强}
{当放射性元素的原子的核外电子具有较高能量时，将发生 $\beta$ 衰变}
{放射性的原子核发生衰变后产生的新核从高能级向低能级跃迁时，辐射出 $\gamma$ 射线}

%% answer: D
%% memo:
\memoanswer{
本题原卷及材料中未提供详解，待补充。
}

\item
%% number: 17
%% paperName: 2007年四川理综第 17 题 6 分
%% typeId: 1
%% type: 单选题
%% chapter: 曲线运动
%% point: 天体运动
%% method: 无
%% score: 6
%% degree: 450
%% duplicateId: 0
%% body:
我国探月的“嫦娥工程”已启动，在不久的将来，我国宇航员将登上月球。假如宇航员在月球上测得摆长为 $l$ 的单摆做小振幅振动的周期为 $T$，将月球视为密度均匀、半径为 $r$ 的球体，则月球的密度为\xzanswer{B}

\fourchoices[answer=B]
{$\frac{\pi l}{3GrT^{2}}$}
{$\frac{3\pi l}{GrT^{2}}$}
{$\frac{16\pi l}{3GrT^{2}}$}
{$\frac{3\pi l}{16GrT^{2}}$}

%% answer: B
%% memo:
\memoanswer{
本题原卷及材料中未提供详解，待补充。
}

\item
%% number: 18
%% paperName: 2007年四川理综第 18 题 6 分
%% typeId: 1
%% type: 单选题
%% chapter: 运动和力
%% point: 动量守恒定律
%% method: 无
%% score: 6
%% degree: 450
%% duplicateId: 0
%% body:
如图所示，弹簧的一端固定在竖直墙上，质量为 $m$ 的光滑弧形槽静止在光滑水平面上，底部与水平面平滑连接，一个质量也为 $m$ 的小球从槽高 $h$ 处开始下滑\xzanswer{C}

\onepicture[w=7.0cm]{figs/18.png}

\fourchoices[answer=C]
{在以后的运动过程中，小球和槽的动量始终守恒}
{在下滑过程中小球和槽之间的相互作用力始终不做功}
{在弹簧反弹后，小球和槽都做速率不变的直线运动}
{被弹簧反弹后，小球和槽的机械能守恒，小球能回到槽高 $h$ 处}

%% answer: C
%% memo:
\memoanswer{
本题原卷及材料中未提供详解，待补充。
}

\item
%% number: 19
%% paperName: 2007年四川理综第 19 题 6 分
%% typeId: 2
%% type: 多选题
%% chapter: 光学
%% point: 光电效应
%% method: 无
%% score: 6
%% degree: 650
%% duplicateId: 0
%% body:
两种单色光 $a$ 和 $b$，$a$ 光照射某金属时有光电子逸出，$b$ 光照射该金属时没有光电子逸出，则\xzanswer{BD}

\fourchoices[answer=BD]
{在真空中，$a$ 光的传播速度较大}
{在水中，$a$ 光的波长较小}
{在真空中，$b$ 光光子能量较大}
{在水中，$b$ 光的折射率较小}

%% answer: BD
%% memo:
\memoanswer{
本题原卷及材料中未提供详解，待补充。
}

\item
%% number: 20
%% paperName: 2007年四川理综第 20 题 6 分
%% typeId: 2
%% type: 多选题
%% chapter: 机械波
%% point: 波与振动图象
%% method: 无
%% score: 6
%% degree: 450
%% duplicateId: 0
%% body:
图甲为一列简谐横波在某一时刻的波形图，图乙为质点 $P$ 以此时刻为计时起点的振动图象。从该时刻起\xzanswer{AC}

\twopicture[numstyle=chinese, w=6.0cm, labelA=2007四川理综20a, labelB=2007四川理综20b]
{figs/20a.png}
{figs/20b.png}

\fourchoices[answer=AC]
{经过 $0.35\Us$ 时，质点 $Q$ 距平衡位置的距离小于质点 $P$ 距平衡位置的距离}
{经过 $0.25\Us$ 时，质点 $Q$ 的加速度大于质点 $P$ 的加速度}
{经过 $0.15\Us$，波沿 $x$ 轴的正方向传播了 $3\Um$}
{经过 $0.1\Us$ 时，质点 $Q$ 的运动方向沿 $y$ 轴正方向}

%% answer: AC
%% memo:
\memoanswer{
本题原卷及材料中未提供详解，待补充。
}

\item
%% number: 21
%% paperName: 2007年四川理综第 21 题 6 分
%% typeId: 1
%% type: 单选题
%% chapter: 磁场
%% point: 带电粒子在磁场中的运动
%% method: 无
%% score: 6
%% degree: 450
%% duplicateId: 0
%% body:
如图所示，长方形 $abcd$ 长 $ad=0.6\Um$，宽 $ab=0.3\Um$，$O$、$e$ 分别是 $ad$、$bc$ 的中点，以 $ad$ 为直径的半圆内有垂直纸面向里的匀强磁场（边界上无磁场），磁感应强度 $B=0.25\UT$。一群不计重力、质量 $m=3\times10^{-7}\Ukg$、电荷量 $q=+2\times10^{-3}\UC$ 的带电粒子以速度 $v=5\times10^{2}\Ums$ 沿垂直 $ad$ 方向且垂直于磁场射入磁场区域\xzanswer{D}

\onepicture[w=3.5cm]{figs/21.png}

\fourchoices[answer=D]
{从 $Od$ 边射入的粒子，出射点全部分布在 $Oa$ 边}
{从 $aO$ 边射入的粒子，出射点全部分布在 $ab$ 边}
{从 $Od$ 边射入的粒子，出射点分布在 $Oa$ 边和 $ab$ 边}
{从 $aO$ 边射入的粒子，出射点分布在 $ab$ 边和 $be$ 边}

%% answer: D
%% memo:
\memoanswer{
本题原卷及材料中未提供详解，待补充。
}

\item
%% number: 22
%% paperName: 2007年四川理综第 22 题 12 分
%% typeId: 4
%% type: 实验题
%% chapter: 电路
%% point: 测电动势与内阻
%% method: 无
%% score: 12
%% degree: 450
%% duplicateId: 0
%% body:
\begin{enumerate}
	\item
	（5 分）在做“研究平抛物体的运动”实验时，除了木板、小球、斜槽、铅笔、图钉之外，下列器材中还需要的是 \tkanswer[1.5cm]{CF}。

	（A）游标卡尺　（B）秒表　（C）坐标纸　（D）天平　（E）弹簧秤　（F）重垂线

	实验中，下列说法正确的是 \tkanswer[1.5cm]{AD}。

	（A）应使小球每次从斜槽上相同的位置自由滑下

	（B）斜槽轨道必须光滑

	（C）斜槽轨道末端可以不水平

	（D）要使描出的轨迹更好地反映真实运动，记录的点应适当多一些

	（E）为了比较准确地描出小球运动的轨迹，应该用一条曲线把所有的点连接起来

	\item
	（12 分）甲同学设计了如图所示的电路测电源电动势 $E$ 及电阻 $R_{1}$ 和 $R_{2}$ 的阻值。实验器材有：待测电源 $E$（不计内阻），待测电阻 $R_{1}$，待测电阻 $R_{2}$，电压表 V（量程为 $1.5\UV$，内阻很大），电阻箱 $R$（$0\sim99.99\UO$），单刀单掷开关 $S_{1}$，单刀双掷开关 $S_{2}$，导线若干。

	\onepicture[w=4.5cm, align=right]{figs/22a.png}

	（1）先测电阻 $R_{1}$ 的阻值。请将甲同学的操作补充完整：闭合 $S_{1}$，将 $S_{2}$ 切换到 $a$，调节电阻箱，读出其示数 $r$ 和对应的电压表示数 $U_{1}$，保持电阻箱示数不变，\tkanswer[2.5cm]{将 $S_{2}$ 切换到 $b$}，读出电压表的示数 $U_{2}$。则电阻 $R_{1}$ 的表达式为 $R_{1}=$ \tkanswer[2.5cm]{$\frac{U_{2}-U_{1}}{U_{1}}r$}。

	（2）甲同学已经测得电阻 $R_{1}=4.8\UO$，继续测电源电动势 $E$ 和电阻 $R_{2}$ 的阻值。该同学的做法是：闭合 $S_{1}$，将 $S_{2}$ 切换到 $a$，多次调节电阻箱，读出多组电阻箱示数 $R$ 和对应的电压表示数 $U$，由测得的数据，绘出了如图所示的 $\frac{1}{U}-\frac{1}{R}$ 图线，则电源电动势 $E=$ \tkanswer[1.2cm]{1.43} $\UV$，电阻 $R_{2}=$ \tkanswer[1.2cm]{1.2} $\UO$。

	\onepicture[w=4.5cm, align=right]{figs/22b.png}

	（3）利用甲同学设计的电路和测得的电阻 $R_{1}$，乙同学测电源电动势 $E$ 和电阻 $R_{2}$ 的阻值的做法是：闭合 $S_{1}$，将 $S_{2}$ 切换到 $b$，多次调节电阻箱，读出多组电阻箱示数 $R$ 和对应的电压表示数 $U$，由测得的数据，绘出了相应的 $\frac{1}{U}-\frac{1}{R+R_{1}}$ 图线，根据图线得到电源电动势 $E$ 和电阻 $R_{2}$。这种做法与甲同学的做法比较，由于电压表测得的数据范围 \tkanswer[1.2cm]{较小}（选填“较大”、“较小”或“相同”），所以 \tkanswer[1.2cm]{甲} 同学的做法更恰当些。
\end{enumerate}

%% answer:
\jdanswer{
\begin{enumerate}
	\item
	CF，AD。

	\item
	（1）将 $S_{2}$ 切换到 $b$，$\frac{U_{2}-U_{1}}{U_{1}}r$；（2）$1.43\UV$（或 $\frac{10}{7}\UV$），$1.2\UO$；（3）较小，甲。
\end{enumerate}
}

%% memo:
\memoanswer{
本题原卷及材料中未提供详解，待补充。
}

\item
%% number: 23
%% paperName: 2007年四川理综第 23 题 16 分
%% typeId: 6
%% type: 计算题
%% chapter: 电磁感应
%% point: 电磁感应受力分析
%% method: 无
%% score: 16
%% degree: 450
%% duplicateId: 0
%% body:
如图所示，$P$、$Q$ 为水平面内平行放置的光滑金属长直导轨，间距为 $L_{1}$，处在竖直向下、磁感应强度大小为 $B_{1}$ 的匀强磁场中。一导体杆 $ef$ 垂直于 $P$、$Q$ 放在导轨上，在外力作用下向左做匀速直线运动。质量为 $m$、每边电阻均为 $r$、边长为 $L_{2}$ 的正方形金属框 $abcd$ 置于竖直平面内，两顶点 $a$、$b$ 通过细导线与导轨相连，磁感应强度大小为 $B_{2}$ 的匀强磁场垂直金属框向里，金属框恰好处于静止状态。不计其余电阻和细导线对 $a$、$b$ 点的作用力。

\begin{enumerate}
	\item
	通过 $ab$ 边的电流 $I_{ab}$ 是多大？
	\item
	导体杆 $ef$ 的运动速度 $v$ 是多大？
\end{enumerate}

\onepicture[w=6.5cm, align=right]{figs/23.png}

%% answer:
\jdanswer{
\begin{enumerate}
	\item
	$I_{ab}=\frac{3mg}{4B_{2}L_{2}}$

	\item
	$v=\frac{3mgr}{4B_{1}B_{2}L_{1}L_{2}}$
\end{enumerate}
}

%% memo:
\memoanswer{
（1）设通过正方形金属框的总电流为 $I$，$ab$ 边的电流为 $I_{ab}$，$dc$ 边的电流为 $I_{dc}$，有 $I_{ab}=\frac{3}{4}I$ ①，$I_{dc}=\frac{1}{4}I$ ②。金属框受重力和安培力，处于静止状态，有 $mg=B_{2}I_{ab}L_{2}+B_{2}I_{dc}L_{2}$ ③。由①②③解得 $I_{ab}=\frac{3mg}{4B_{2}L_{2}}$ ④。

（2）由（1）可得 $I=\frac{mg}{B_{2}L_{2}}$ ⑤。设导体杆切割磁感线产生的电动势为 $E$，有 $E=B_{1}L_{1}v$ ⑥。设 $ad$、$dc$、$cb$ 三边电阻串联后与 $ab$ 边电阻并联的总电阻为 $R$，则 $R=\frac{3}{4}r$ ⑦。根据闭合电路欧姆定律，有 $I=E/R$ ⑧。由⑤～⑧解得 $v=\frac{3mgr}{4B_{1}B_{2}L_{1}L_{2}}$ ⑨。
}

\item
%% number: 24
%% paperName: 2007年四川理综第 24 题 19 分
%% typeId: 6
%% type: 计算题
%% chapter: 电场
%% point: 带电粒子的直线运动
%% method: 无
%% score: 19
%% degree: 450
%% duplicateId: 0
%% body:
如图所示，一根长 $L=1.5\Um$ 的光滑绝缘细直杆 $MN$，竖直固定在场强为 $E=1.0\times10^{5}\UNC$、与水平方向成 $\theta=30^{\circ}$ 角的倾斜向上的匀强电场中。杆的下端 $M$ 固定一个带电小球 A，电荷量 $Q=+4.5\times10^{-6}\UC$；另一带电小球 B 穿在杆上可自由滑动，电荷量 $q=+1.0\times10^{-6}\UC$，质量 $m=1.0\times10^{-2}\Ukg$。现将小球 B 从杆的上端 N 静止释放，小球 B 开始运动。（静电力常量 $k=9.0\times10^{9}\UNmqCq$，取 $g=10\Umsq$）

\begin{enumerate}
	\item
	小球 B 开始运动时的加速度为多大？
	\item
	小球 B 的速度最大时，距 M 端的高度 $h_{1}$ 为多大？
	\item
	小球 B 从 N 端运动到距 M 端的高度 $h_{2}=0.61\Um$ 时，速度为 $v=1.0\Ums$，求此过程中小球 B 的电势能改变了多少？
\end{enumerate}

\onepicture[w=3.5cm, align=right]{figs/24.png}

%% answer:
\jdanswer{
\begin{enumerate}
	\item
	$a=3.2\Umsq$

	\item
	$h_{1}=0.9\Um$

	\item
	$\Delta E_{P}=8.2\times10^{-2}\UJ$
\end{enumerate}
}

%% memo:
\memoanswer{
（1）开始运动时小球 B 受重力、库仑力、杆的弹力和电场力，沿杆方向运动，由牛顿第二定律得 $mg-k\frac{Qq}{L^{2}}-qE\sin\theta=ma$ ①。解得 $a=g-\frac{kQq}{L^{2}m}-\frac{qE\sin\theta}{m}$ ②。代入数据解得 $a=3.2\Umsq$ ③。

（2）小球 B 速度最大时合力为零，即 $\frac{kQq}{h_{1}^{2}}+qE\sin\theta=mg$ ④。解得 $h_{1}=\sqrt{\frac{kQq}{mg-qE\sin\theta}}$ ⑤。代入数据解得 $h_{1}=0.9\Um$ ⑥。

（3）小球 B 从开始运动到速度为 $v$ 的过程中，设重力做功为 $W_{1}$，电场力做功为 $W_{2}$，库仑力做功为 $W_{3}$，根据动能定理有 $W_{1}+W_{2}+W_{3}=\frac{1}{2}mv^{2}$ ⑦，$W_{1}=mg(L-h_{2})$ ⑧，$W_{2}=-qE(L-h_{2})\sin\theta$ ⑨。解得 $W_{3}=\frac{1}{2}mv^{2}-mg(L-h_{2})+qE(L-h_{2})\sin\theta$ ⑩。设小球的电势能改变了 $\Delta E_{p}$，则 $\Delta E_{p}=-(W_{2}+W_{3})$（11），$\Delta E_{p}=mg(L-h_{2})-\frac{1}{2}mv^{2}$（12），$\Delta E_{p}=8.2\times10^{-2}\UJ$（13）。
}

\item
%% number: 25
%% paperName: 2007年四川理综第 25 题 20 分
%% typeId: 6
%% type: 计算题
%% chapter: 运动和力
%% point: 动量和能量
%% method: 无
%% score: 20
%% degree: 250
%% duplicateId: 0
%% body:
目前，滑板运动受到青少年的追捧。如图是某滑板运动员在一次表演时的一部分赛道在竖直平面内的示意图。赛道光滑，$FGI$ 为圆弧赛道，半径 $R=6.5\Um$，$C$ 为最低点并与水平赛道 $BC$ 位于同一水平面，$KA$、$DE$ 平台的高度都为 $h=1.8\Um$。$B$、$C$、$F$ 处平滑连接。滑板 $a$ 和 $b$ 的质量均为 $m$，$m=5\Ukg$，运动员质量为 $M$，$M=45\Ukg$。

表演开始，运动员站在滑板 $b$ 上。先让滑板 $a$ 从 $A$ 点静止下滑，$t_{1}=0.1\Us$ 后再与 $b$ 板一起从 $A$ 点静止下滑。滑上 $BC$ 赛道后，运动员从 $b$ 板跳到同方向运动的 $a$ 板上，在空中运动的时间 $t_{2}=0.6\Us$（水平方向是匀速运动）。运动员与 $a$ 板一起沿 $CD$ 赛道上滑后冲出赛道，落在 $EF$ 赛道的 $P$ 点，沿赛道滑行，经过 $G$ 点时，运动员受到的支持力 $N=742.5\UN$。（滑板和运动员的所有运动都在同一竖直平面内，计算时滑板和运动员都看作质点，取 $g=10\Umsq$）

\begin{enumerate}
	\item
	滑到 G 点时，运动员的速度是多大？
	\item
	运动员跳上滑板 $a$ 后，在 $BC$ 赛道上与滑板 $a$ 共同运动的速度是多大？
	\item
	从表演开始到运动员滑至 I 的过程中，系统的机械能改变了多少？
\end{enumerate}

\onepicture[w=7.5cm, align=right]{figs/25.png}

%% answer:
\jdanswer{
\begin{enumerate}
	\item
	$6.5\Ums$

	\item
	$6.9\Ums$

	\item
	$88.75\UJ$
\end{enumerate}
}

%% memo:
\memoanswer{
（1）在 G 点，运动员和滑板一起做圆周运动，设向心加速度为 $a_{\text{向}}$，速度为 $v_{G}$，运动员受到重力 $Mg$、滑板对运动员的支持力 $N$ 的作用，则 $N-Mg=Ma_{\text{向}}$ ①，$a_{\text{向}}=\frac{v_{G}^{2}}{R}$ ②，$N-Mg=M\frac{v_{G}^{2}}{R}$ ③，$v_{G}=\sqrt{\frac{R(N-Mg)}{M}}$ ④，$v_{G}=6.5\Ums$ ⑤。

（2）设滑板 $a$ 由 $A$ 点静止下滑到 $BC$ 赛道后速度为 $v_{1}$，由机械能守恒定律有 $mgh=\frac{1}{2}mv_{1}^{2}$ ⑥，$v_{1}=\sqrt{2gh}$ ⑦。运动员与滑板 $b$ 一起由 $A$ 点静止下滑到 $BC$ 赛道后，速度也为 $v_{1}$。运动员由滑板 $b$ 跳到滑板 $a$，设蹬离滑板 $b$ 时的水平速度为 $v_{2}$，在空中飞行的水平位移为 $s$，则 $s=v_{2}t_{2}$ ⑧。设起跳时滑板 $a$ 与滑板 $b$ 的水平距离为 $s_{0}$，则 $s_{0}=v_{1}t_{1}$ ⑨。设滑板 $a$ 在 $t_{2}$ 时间内的位移为 $s_{1}$，则 $s_{1}=v_{1}t_{2}$ ⑩，$s=s_{0}+s_{1}$（11），即 $v_{2}t_{2}=v_{1}(t_{1}+t_{2})$（12）。运动员落到滑板 $a$ 后，与滑板 $a$ 共同运动的速度为 $v$，由动量守恒定律有 $mv_{1}+Mv_{2}=(m+M)v$（13）。由以上方程可解出 $v=\frac{mt_{2}+M(t_{1}+t_{2})}{(M+m)t_{2}}\sqrt{2gh}$（14）。代入数据，解得 $v=6.9\Ums$（15）。

（3）设运动员离开滑板 $b$ 后，滑板 $b$ 的速度为 $v_{3}$，有 $Mv_{2}+mv_{3}=(M+m)v_{1}$（16）。可算出 $v_{3}=-3\Ums$，有 $|v_{3}|=3\Ums<v_{1}=6\Ums$，$b$ 板将在两个平台之间来回运动，机械能不变。系统的机械能改变为 $\Delta E=\frac{1}{2}(M+m)v_{G}^{2}+\frac{1}{2}mv_{3}^{2}-(m+m+M)gh$（17），$\Delta E=88.75\UJ$（18）。
}

\end{enumerate}

\end{document}
